\documentclass[12pt,AutoFakeBold]{ctexart}
\usepackage[a4paper,margin=2.5cm]{geometry}
\usepackage{array}
\usepackage{longtable}
\usepackage{booktabs}
\usepackage{xcolor}          % 支持 HTML 十六进制颜色
\usepackage{colortbl}        % \rowcolor
% 关键：arydshln 必须放在 array / longtable / colortbl 之后加载
\usepackage{arydshln}
\setlength{\dashlinedash}{1.5pt}   % 虚线线段长度（0.5pt 太短，几乎看不见）
\setlength{\dashlinegap}{2pt}      % 虚线间距
% 黑体没有真粗体，开启伪粗体可消除 Font shape undefined 警告
\xeCJKsetup{AutoFakeBold=true}

\renewcommand{\arraystretch}{1.35}
\setlength{\tabcolsep}{4pt}

\newcommand{\levelone}[1]{%
  \par\addvspace{1.5em}%
  \noindent{\heiti\fontsize{15pt}{18pt}\selectfont\color[HTML]{1F3A5F}#1}\par
  \addvspace{0.8em}%
}

\begin{document}

\begin{center}
\textbf{{\heiti\fontsize{22pt}{28pt}\selectfont 2026全国大学生数学建模竞赛}\\[8pt]}
\textbf{{\heiti\fontsize{22pt}{28pt}\selectfont AI工具使用详情说明}}
\end{center}

\levelone{一、所用AI工具名称、版本或型号}
\begingroup\small
\begin{longtable}{c>{\centering\arraybackslash}p{2.8cm}>{\centering\arraybackslash}p{3.5cm}>{\raggedright\arraybackslash}p{8cm}}
\toprule
\rowcolor[HTML]{E8EDF3}
\multicolumn{1}{c}{序号} & \multicolumn{1}{c}{AI工具名称} & \multicolumn{1}{c}{版本/型号} & \multicolumn{1}{c}{主要用途} \\
\midrule
示例 & ChatGPT & GPT-4o & 模型思路参考、代码调试、文字润色 \\
\hdashline
1 & & & \\
\hdashline
2 & & & \\
\hdashline
3 & & & \\
\hdashline
4 & & & \\
\bottomrule
\end{longtable}
\endgroup

\levelone{二、具体使用目的和环节}
\begingroup\small
\begin{longtable}{c>{\raggedright\arraybackslash}p{5cm}c>{\raggedright\arraybackslash}p{4.65cm}>{\raggedright\arraybackslash}p{2.6cm}}
\toprule
\rowcolor[HTML]{E8EDF3}
\multicolumn{1}{c}{序号} & \multicolumn{1}{c}{论文写作环节} & \multicolumn{1}{c}{是否使用} & \multicolumn{1}{c}{使用目的} & \multicolumn{1}{c}{使用工具} \\
\midrule
1 & 赛题理解与问题分析 & 是 & （如：辅助拆解题目要求，识别关键约束条件） & \\
\hdashline
2 & 模型假设与符号定义 & 是 & （如：参考常见假设类型，本队筛选后确定） & \\
\hdashline
3 & 模型建立与算法设计 & 是 & （如：提供算法选型建议，核心建模由本队主导） & \\
\hdashline
4 & 模型求解与编程实现 & 是 & （如：辅助编写求解代码，调试报错，本队验证） & \\
\hdashline
5 & 结果分析与模型检验 & 否 & & \\
\hdashline
6 & 论文撰写与文字润色 & 是 & （如：优化文字表达，调整语句通顺度） & \\
\hdashline
7 & 其他辅助环节（文献/数据/图表等） & 是 & （如：辅助整理文献格式，数据清洗） & \\
\bottomrule
\end{longtable}
\endgroup

\levelone{三、主要提示方式与使用过程说明}
\subsection*{（一）主要提示方式}
\begingroup\small
\begin{longtable}{c>{\raggedright\arraybackslash}p{5cm}c>{\raggedright\arraybackslash}p{7.4cm}}
\toprule
\rowcolor[HTML]{E8EDF3}
\multicolumn{1}{c}{序号} & \multicolumn{1}{c}{提示方式} & \multicolumn{1}{c}{是否使用} & \multicolumn{1}{c}{简要说明} \\
\midrule
1 & 网页对话框交互 & 是/否 & （如：用于问题咨询、思路讨论、文字润色） \\
\hdashline
2 & 代码编辑器内嵌AI & 是/否 & （如：用于代码补全、调试、重构） \\
\hdashline
3 & 上传文件或数据对话 & 是/否 & （如：上传赛题数据或论文草稿进行分析） \\
\hdashline
4 & AI智能体工作流（多步自动执行） & 是/否 & （如：使用自动化Agent完成数据处理流程） \\
\hdashline
5 & 其他方式 & 是/否 & （请说明） \\
\bottomrule
\end{longtable}
\endgroup

\subsection*{（二）典型交互示例}
\noindent 【填写提示】选取竞赛过程中2至3次与AI的关键交互进行记录。选择标准：①对建模思路或结果有重要影响的交互；②能体现本队如何使用和核验AI输出的交互，提示词应完整记录，AI回复核心内容，本队处理方式要写清楚是直接采纳、修改后采纳还是未采纳，以及具体修改了什么。
\medskip

\noindent\textbf{示例1：}
\begingroup\small
\begin{longtable}{>{\raggedright\arraybackslash}m{2cm}>{\raggedright\arraybackslash}m{13cm}}
\toprule
\rowcolor[HTML]{E8EDF3}
\multicolumn{1}{c}{项目} & \multicolumn{1}{c}{填写内容} \\
\midrule
对应环节 & （如：模型建立与算法设计） \\
\hdashline
使用工具及版本 & （如：ChatGPT GPT-4o） \\
\hdashline
本队提示词 & （完整记录向AI提出的问题或指令，可附截图。注意：不要粘贴API Key、账号密码等敏感信息） \\
\hdashline
AI回复内容 & （记录AI的回复核心内容，关键部分可原文引用） \\
\hdashline
本队处理方式 & 直接采纳 / 修改后采纳 / 未采纳 / （说明：修改了哪些内容？为什么修改？未采纳的原因是什么？） \\
\hdashline
人工核验方式 & （说明如何确认AI回复的正确性：如独立推导验证、代码独立复现、数据交叉核对、文献查证、队员讨论确认等） \\
\bottomrule
\end{longtable}
\endgroup
\medskip

\noindent\textbf{示例2：}
\begingroup\small
\begin{longtable}{>{\raggedright\arraybackslash}m{2cm}>{\raggedright\arraybackslash}m{13cm}}
\toprule
\rowcolor[HTML]{E8EDF3}
\multicolumn{1}{c}{项目} & \multicolumn{1}{c}{填写内容} \\
\midrule
对应环节 & （如：模型建立与算法设计） \\
\hdashline
使用工具及版本 & （如：ChatGPT GPT-4o） \\
\hdashline
本队提示词 & （完整记录向AI提出的问题或指令，可附截图。注意：不要粘贴API Key、账号密码等敏感信息） \\
\hdashline
AI回复内容 & （记录AI的回复核心内容，关键部分可原文引用） \\
\hdashline
本队处理方式 & 直接采纳 / 修改后采纳 / 未采纳 / （说明：修改了哪些内容？为什么修改？未采纳的原因是什么？） \\
\hdashline
人工核验方式 & （说明如何确认AI回复的正确性：如独立推导验证、代码独立复现、数据交叉核对、文献查证、队员讨论确认等） \\
\bottomrule
\end{longtable}
\endgroup

\levelone{四、对AI输出的采纳、人工修改和核验的主要情况}
\noindent 【填写提示】对AI输出内容（语言润色除外）的采纳、人工修改和核验情况进行总体说明。注意：①语言润色类内容无需逐项说明；②核心建模与分析须由本队主导，AI仅作辅助；③“采纳与修改情况”栏要写清楚是部分修改后采纳，并简述修改内容；④“人工核验方式”栏要写清楚本队如何验证AI输出的正确性，不能只写“已核验”。
\medskip

\begingroup\small
\begin{longtable}{c>{\raggedright\arraybackslash}p{3.7cm}>{\raggedright\arraybackslash}p{4.6cm}>{\raggedright\arraybackslash}p{6cm}}
\toprule
\rowcolor[HTML]{E8EDF3}
\multicolumn{1}{c}{序号} & \multicolumn{1}{c}{AI输出内容类别} & \multicolumn{1}{c}{采纳与修改情况} & \multicolumn{1}{c}{人工核验方式} \\
\midrule
1 & 建模思路与方法建议 & （如：AI提供3种思路，本队选择第2种并改进） & （如：本队独立推导验证，对比文献确认可行性） \\
\hdashline
2 & 公式推导与理论参考 & & \\
\hdashline
3 & 代码编写与调试 & & \\
\hdashline
4 & 其他内容 & & \\
\bottomrule
\end{longtable}
\endgroup
\medskip

\noindent\textbf{核心环节人工主导确认：}
\begingroup\small
\begin{longtable}{c>{\raggedright\arraybackslash}p{3.5cm}c>{\raggedright\arraybackslash}p{8.5cm}}
\toprule
\rowcolor[HTML]{E8EDF3}
\multicolumn{1}{c}{序号} & \multicolumn{1}{c}{核心环节} & \multicolumn{1}{c}{是否本队主导} & \multicolumn{1}{c}{本队贡献说明} \\
\midrule
1 & 模型结构与创新点 & 是 & \\
\hdashline
2 & 公式推导与求解步骤 & 是 & \\
\hdashline
3 & 程序逻辑与参数设置 & 是 & \\
\hdashline
4 & 结果分析与论文核心论述 & 是 & \\
\hdashline
5 & 论文撰写与图表制作等 & 是 & \\
\bottomrule
\end{longtable}
\endgroup

\medskip
\noindent 本队确认：以上AI工具使用情况真实完整，无隐瞒、无虚假陈述。核心建模与分析由本队主导完成，所有AI输出内容（语言润色除外）均经过人工审查与核实。如有不实，本队愿承担相应责任。

\end{document}
